\section*{СПИСОК СОКРАЩЕНИЙ И ОБОЗНАЧЕНИЙ}
\addcontentsline{toc}{section}{СПИСОК СОКРАЩЕНИЙ И ОБОЗНАЧЕНИЙ}

\begin{longtable}{p{4cm} p{11cm}}
\toprule
\textbf{Сокращение} & \textbf{Расшифровка} \\
\midrule
\endfirsthead
\multicolumn{2}{c}{\tablename\ \thetable\ --- продолжение} \\
\toprule
\textbf{Сокращение} & \textbf{Расшифровка} \\
\midrule
\endhead
\bottomrule
\endlastfoot

AST    & Abstract Syntax Tree --- абстрактное синтаксическое дерево \\
CFG    & Control Flow Graph --- граф управляющего потока \\
CUDA   & Compute Unified Device Architecture --- унифицированная архитектура вычислений \\
DFT    & Discrete Fourier Transform --- дискретное преобразование Фурье \\
DIT    & Decimation In Time --- прореживание по времени \\
ELF    & Executable and Linkable Format --- формат исполняемых и компонуемых файлов \\
FFT    & Fast Fourier Transform --- быстрое преобразование Фурье \\
GELU   & Gaussian Error Linear Unit --- нелинейность на основе функции ошибки \\
GOT    & Global Offset Table --- таблица глобальных смещений \\
GPU    & Graphics Processing Unit --- графический процессор \\
IEEE   & Institute of Electrical and Electronics Engineers \\
IPDOM  & Immediate Post-Dominator --- непосредственный постдоминатор \\
IR     & Intermediate Representation --- промежуточное представление \\
ISA    & Instruction Set Architecture --- архитектура системы команд \\
JIT    & Just-In-Time --- динамическая (just-in-time) компиляция \\
NFA    & Nondeterministic Finite Automaton --- недетерминированный конечный автомат \\
NTP    & Network Time Protocol --- протокол синхронизации времени \\
OCP    & Open-Closed Principle --- принцип открытости-закрытости \\
PC     & Program Counter --- программный счётчик \\
PDOM   & Post-DOMinator --- постдоминатор \\
PLT    & Procedure Linkage Table --- таблица связей процедур \\
PTX    & Parallel Thread Execution --- параллельное выполнение потоков (ISA NVIDIA) \\
SASS   & Shader ASSembly --- машинный язык GPU NVIDIA \\
SIMT   & Single Instruction, Multiple Threads --- одна инструкция, несколько потоков \\
СЛАУ   & Система линейных алгебраических уравнений \\
SM     & Streaming Multiprocessor --- потоковый мультипроцессор \\
\midrule
\multicolumn{2}{l}{\textbf{Математические обозначения}} \\
\midrule
$\mathbb{B}^{32}$ & пространство двоичных векторов размера 32 \\
$\lfloor x \rfloor$ & целая часть (пол) числа $x$ \\
$\lceil x \rceil$  & наименьшее целое, не меньшее $x$ (потолок) \\
$|S|$              & мощность множества $S$ \\
$\popcount(x)$     & число единичных бит в $x$ \\
$\mathrm{round}_\xi$ & округление в режиме $\xi$ \\
$\mathrm{active}(M)$ & множество индексов активных потоков при маске $M$ \\
$\eta_{\mathrm{branch}}$ & метрика эффективности ветвления \\
$\kappa$           & коэффициент слияния глобальных транзакций \\
$C_{\mathrm{bank}}$ & число конфликтов банков разделяемой памяти \\
$T_{\mathrm{actual}}, T_{\mathrm{ideal}}$ & фактическое и идеальное число L1-транзакций \\
$W = \langle T, M, \mathrm{PC}, R, S \rangle$ & формальное представление варпа \\
\end{longtable}

\newpage
